\appendix
\section{Full Reviewer Reports}
\label{app:reviews}

We reproduce below the complete text of the editor's decision letter,
the Associate Editor's report, and the two referee reports, exactly
as received.
These are included for reference so that the mapping between
individual comments and revision actions
(Table~\ref{tab:themes}, Table~\ref{tab:problem_theme_map})
can be verified.


% =========================================================
\subsection{Editor Decision Letter (Po-Ling Loh)}
\label{app:editor}
% =========================================================

\begin{quote}
\small
Dear Mario Beraha,

Thank you for submitting your paper ``Online activity prediction via
generalized Indian buffet process models'' for possible publication
in Annals of Applied Statistics.
It has now been carefully reviewed and my decision is:
\textbf{Revision required.}

Thank you for submitting your paper for possible publication in
Annals of Applied Statistics. After careful reviews, our decision is:
\textbf{Major revision required.}

The paper has been reviewed by an Associate Editor and two referees.

[\emph{AE report reproduced in Section~\ref{app:ae} below.}]

Given these concerns, we cannot accept your manuscript for publication
in its current form. Meanwhile, we are open to considering a
thoroughly revised version that fully addresses all questions raised,
with the understanding that there is no guarantee for the ultimate
outcome of the new version at this moment.

Should you choose to revise the manuscript, please keep the following
in mind.
\begin{itemize}[leftmargin=*]
  \item It is strongly recommended to include a dedicated section,
        typically Section~2, to introduce the motivating application
        and dataset, and to clearly explain why it is important to
        address the applied questions at hand. Then in the real-world
        data analysis section, please provide sufficient details on
        the applied findings and their implications for the field.
  \item It is strongly recommended to use the template file provided
        by the journal. Please note that most published papers will
        not exceed 20~pages (about 500 words per page, with figures
        typically taking 1/3~page and displayed equations 30~words
        each), anything longer requiring unusually compelling subject
        matter.
  \item Printing figures in color adds significantly to the production
        cost of the journal. While color may be used in the online
        publication, we will use color in the printed version only
        when essential to the display. Please use dashed/dotted lines
        or symbols where possible and avoid referring to colors in the
        text and the figure caption.
  \item It is strongly recommended to submit the revision within
        4~months of the decision date. A revision not submitted within
        6~months will be automatically withdrawn.
  \item Please prepare a detailed point-by-point response letter to
        the Editor, the Associate Editor, and each individual referee,
        all in a single file.
\end{itemize}

Sincerely,\\
Po-Ling Loh\\
Editor, The Annals of Applied Statistics
\end{quote}


% =========================================================
\subsection{Associate Editor Report}
\label{app:ae}
% =========================================================

\begin{quote}
\small
Both reviewers agree that this paper makes a nice methodological
contribution, and that the motivation for the paper is compelling.
However, this paper is not yet suitable for AoAS, given the applied
and real world focus of the journal. The practical impact of the
methodology is unclear, as the simulations do not clearly show this
method is better, and the data analysis is not very thorough and it
is unclear what the significance of this dataset is. Perhaps the
authors should focus more on the data analysis and the downstream
impact of the methods they propose.

This methodology is highly related to Camerlenghi et al (2022). This
manuscript should more clearly articulate its original contributions
and how they extend beyond existing work.

There are also multiple typos and confusing points as noted by the
reviewers, that should be corrected.
\end{quote}


% =========================================================
\subsection{Reviewer~1 Report}
\label{app:r1}
% =========================================================

\begin{quote}
\small
\textbf{General comment.}
The authors introduce a novel Bayesian nonparametric framework for
the prediction of count data. The paper is motivated by a clear and
relevant application: predicting new users, total triggers, and the
sufficient length of the follow-up period in the design of A/B tests
(i.e., an online marketing methodology aimed at comparing the
performance of two versions of a product or service). Building upon
the literature on the Indian buffet process and its generalizations,
the authors' approach yields closed-form posterior and predictive
quantities that are both theoretically appealing and computationally
efficient. In particular, parameter estimation and prediction under
the proposed models do not require computationally expensive Markov
chain Monte Carlo methods, as is often the case in Bayesian
nonparametric models. This results in a theoretically elegant and
practically flexible framework that can be applied to large datasets
thanks to its computational efficiency. I thoroughly enjoyed reading
this paper. It is well written and carefully organized, with strong
applied motivation, solid methodological development, and convincing
validation through extensive simulations and real-world experiments.
However, I believe that some adjustments are needed to make the paper
fully ready for publication; those are detailed below.

\medskip
\textbf{Point 1.}
While the paper is overall well written, I found Section~3.4 the most
difficult to read and somewhat confusing. Most of the confusion likely
stems from the fact that $D_M$, defined here as the number of
follow-up days needed to see $M$ new users, is set equal to
$F'_{M - N_{D_0}}$, i.e., the day when the $(M - N_{D_0})$-th new
user triggered, which does not seem to align with the earlier
definition. Is the term ``new users'' used interchangeably with
``unobserved users'' in this section? Do both terms refer to users
not observed (i.e., not triggered) during the pilot study of length
$D_0$?

Moreover, while the measure $F'$ is clearly related to the measure
$Z^{FT}$ in Eq.~(4), the presence of Theorem~3.8 in this section
suggests that the reader would benefit from further clarification
here either on the distributional assumptions for $F'$ or through an
explicit reference describing its relationship with $Z^{FT}$.
Without such clarification, the statement of Theorem~3.8 appears
somewhat vague.

\medskip
\textbf{Point 2.}
p.3, Eq.~(1): the second sum should run from $D_0 + 1$ to
$D_0 + D_1$, right?

\medskip
\textbf{Point 3.}
p.4, ll.8--10: the sentence starting from ``this assumption
clearly\ldots'' sounds too strong, or at least imprecise. At this
point in the paper, no distributional assumption on the observations
$A_{d,n}$ has been specified. Why is it obvious that all users out of
the fixed total number $M$ will eventually trigger as $D_1 \to \infty$?

\medskip
\textbf{Point 4.}
In Definition~3.1, the statement
$A_{d,n} \mid \mu \stackrel{\mathrm{iid}}{\sim} G(\cdot \mid \theta_n)$
would likely be clearer if written as
$A_{d,n} \mid \mu \stackrel{\mathrm{ind}}{\sim} G(\cdot \mid \theta_n)$,
since the distribution on the right-hand side depends on $n$.

\medskip
\textbf{Point 5.}
On p.5, Eq.~(5), there appears to be an $s$ in place of a $\theta$.
Moreover, the subsequent phrasing ``which puts to zero all values
greater than $D$'' sounds too informal; consider rephrasing it more
formally, for example, as ``which assigns zero probability to all
values exceeding $D$.''

\medskip
\textbf{Point 6.}
On p.5/6, I suggest mentioning that $P_0$ is a non-atomic
distribution, either in the last paragraph of p.5 or within
Definition~3.2.

\medskip
\textbf{Point 7.}
At the end of Definition~3.2, is $\tilde\mu$ supposed to be just $\mu$?

\medskip
\textbf{Point 8.}
In Theorem~3.4, was $\pi_{\mathcal{G}}$ previously defined?

\medskip
\textbf{Point 9.}
When specific theorems and the algorithm are referred to within the
text, the words ``theorem'' and ``algorithm'' are sometimes
capitalized and sometimes not.

\medskip
\textbf{Point 10.}
In Proposition~3.6, the sum defining $U_{D_0}^{(D_1)}$ should start
from $D_0 + 1$.

\medskip
\textbf{Point 11.}
p.9: there is a typo and the word ``users'' is written as ``usres.''
\end{quote}


% =========================================================
\subsection{Reviewer~2 Report}
\label{app:r2}
% =========================================================

\begin{quote}
\small
The paper proposes an applied-methods framework to forecast user
engagement in online A/B tests: (i)~number of new/distinct users,
(ii)~total triggers including repeats, and (iii)~days to reach target
$M$ users. It builds on stable-beta scaled process priors to model
heavy-tailed user propensities and obtains closed-form
posterior/predictive distributions that make the approach fast for
large platforms.

\medskip
\textbf{Major concern (a).}
Although the context and the theoretical results are tailored to the
specific application in A/B testing, the backbone of the proposed
framework is essentially the same as in Camerlenghi et al (2022),
that is, the scaled-process priors, in particular the stable-Beta
scaled process (SB-SP), used to overcome the rigidity of standard CRM
priors for ``unseen features.'' It's clear that the overlap is mostly
at the level of the prior construction and its posterior-predictive
flexibility; the differences lie in inferential targets
(unseen-feature estimation vs.\ multi-horizon forecasting and stopping
rules) and the emphasis on deployable procedures and calibration for
A/B experimentation workflows. This is my conclusion after reading
both manuscripts; however, this manuscript should more clearly
articulate its original contributions and how they extend beyond
existing work.

\medskip
\textbf{Major concern (b).}
The manuscript is interesting, extensive, and technically careful, but
my major concern is that its center of gravity is on probabilistic
development: multiple likelihood types, several inferential targets,
and a rich catalogue of distributional results. Therefore, it reads
more like a methodological contribution adapted to an applied venue
than a paper conceived with a truly applied motivation. The data
analysis feels thin: a couple of datasets only superficially
investigated, sparse diagnostics, and little evidence that the method
has a practical impact and changes decisions in A/B testing (e.g.,
earlier stopping, fewer underpowered launches, calibrated risk at
given service-level targets).

\medskip
\textbf{Major concern (c).}
The paper's general objective is unclear. The introduction oscillates
between (i)~improving causal effect estimation and (ii)~improving
experimental sampling/design, but it never nails down a single primary
goal, the estimand, or the decision the method is meant to support.
If the contribution is design, the paper must say what decision is
being optimized and under which constraints, e.g., ``choose $D_0$
(pilot length) and/or $n_0$ (pilot size) to minimize expected total
time/cost to a pre-specified precision on the ATE.'' If the true goal
is causal estimation rather than design, the authors should define the
estimand and demonstrate precision/coverage gains relative to standard
estimators, treating design results as ancillary. Moreover, the data
analysis should demonstrate the method's impact and its practical
advantage over competing methods in the context of A/B testing.

\medskip
\textbf{Major concern (d).}
Theoretical results are shown without comments about their
consequences for the application. For example, after Corollary~3.5,
it is mentioned that the predictive distribution of quantities that
can be expressed as functionals of $\mu'$ will be the same under the
Be-SSP and TG-SSP models. However, the practical consequences are not
discussed. In Proposition~3.7, the Bayesian estimator for
$T_{D_0}^{(D_1)}$ has a potentially interesting mixture form, e.g.,
in terms of rate ratios, but this is not commented upon.

\medskip
\textbf{Minor points and editorial issues.}

\medskip
\textbf{(a)} p.3: ``One of the main challenges\ldots crucial for
accurately predicting user engagement.'' The motivation for the paper
would be clearer if those features were visualized in the data.

\textbf{(b)} p.3: ``We can broadly categorize\ldots linear
programming approach.'' It is not clear what the difference is between
the proposed method and the approaches mentioned here, and why these
approaches are limited.

\textbf{(c)} p.4: ``Ionita-Laza et al.\ (2009); Richardson et al.\
(2022) assume that the total number of users, say $M$, is fixed and
known.'' Isn't it always the case, though, that the number of users
is known, in A/B testing? Isn't there a cost to add days and subjects?

\textbf{(d)} p.4: ``user activity often exhibits a power law growth.''
Please, add a reference or a visualization.

\textbf{(e)} p.4: ``Moreover, in some situations, focusing only on
the first trigger times $F_n$ might be preferable as we demonstrate
in our experiments.'' It looks like $F_n$ can be derived from coarse
data, so this point is not clear.

\textbf{(f)} The relevance of Definition~3.1 is clearly not related
to the application. Moreover, it introduces concepts that may not be
familiar to the reader, like the score distribution, which has not
been previously defined.

\textbf{(g)} p.5: ``Parametrizing the success probability of the
negative binomial distribution\ldots infinitesimally small values.''
This sentence needs clarification.

\textbf{(h)} p.5: ``the choice of the prior distribution for $\mu$
is crucial for capturing the complex dynamics and heterogeneity of
user behavior.'' This sentence needs clarification.

\textbf{(i)} Section~5.2 includes some comparisons with other
methods. However, Figure~6 shows that across values of $\tau$, the
proposed Be-SSP has predictive performance comparable to that of the
Indian buffet process (IBP). Because the scaled-stable-beta-Bernoulli
process (SSP) is not plotted, the claimed comparability to SSP cannot
be evaluated. As a result, the empirical advantage of the proposed
method is not clear in this setting.

\textbf{(j)} The real-data analysis would benefit from a fuller
description of the dataset. Even if the data are proprietary, the
paper should provide high-level information about the types of
experiments conducted, the goals of the analysis, and why the proposed
model is suitable for those aims. It is also unclear how the results
would inform company decision-making, particularly relative to
alternative modeling proposals.

\textbf{(k)} If it were for me, I would move most of the Simulations
in the Supplementary Material, and focus the manuscript on the impact
on the real data analysis.

\textbf{(l)} The authors correctly ``highlight how the choice of the
prior distribution for $\mu$ plays a crucial role for the quality of
the posterior and predictive distributions of interest.'' I would also
add the choice of the score distribution. However, the manuscript does
not thoroughly discuss how to assess model fit and select the optimal
score distribution for a given problem.

\medskip
\textbf{Typos and inaccuracies.}
The paper contains numerous typos and inaccuracies. Some
cross-references to equations and citations are unresolved (??). The
paper needs an accurate reading. I am listing several instances here;
however, the issues are present throughout the article.
\begin{itemize}[leftmargin=*]
  \item p.3: If $F_n = \min_d \mathbb{I}[A_{d,n} > 0]$ then the set
        of tuples does not need to include $F_n$. The data are simply
        represented by $A_{1:D,n}$.
  \item p.5, Section~3.2: ``$\Delta^{-\alpha} \sim \mathcal{E}(1)$'';
        $\mathcal{E}(1)$ is not defined.
  \item p.7, Corollary~3.5: What is
        $\tilde\Delta_{1,h}^{-\alpha}$ in (9)? It is not defined in
        the Corollary.
  \item p.11: ``our strategy is more scalable as it does require
        performing Markov chain Monte Carlo simulations'' $\to$
        ``does \emph{not} require.''
  \item p.16: ``Moreover, the length of the intervals increases with
        parameter tau of the data generating process'' $\to$ the
        parameter $\tau$.
\end{itemize}
\end{quote}
